%% kapa: exam=1 round=32 year=2021 period=1 form=A subject=economics subject_ko=경제학원론 seq=39 number=79
%% source: Q-Net 감정평가사 기출문제 (cst00309) artlSeq=5208397 file=extracted/1008/5208397/A형/2021년도 제32회 감정평가사 1차 1교시 A형.pdf
%% conversion: kordoc 4.24.51 -> tools/gen1.py (auto); review=proofread-sonnet; flags=none
\begin{question}{79}
예상보다 높은 인플레이션이 발생할 경우 나타나는 효과에 관한 설명으로 옳지 않은 것은?
\begin{choices}
  \item 누진세 체계 하에서 정부의 조세수입은 감소한다.
  \item 채무자는 이익을 보지만 채권자는 손해를 보게 된다.
  \item 고정된 화폐소득을 얻는 봉급생활자는 불리해진다.
  \item 명목 국민소득이 증가한다.
  \item 화폐의 구매력이 감소한다.
\end{choices}
\end{question}
\answer{1}
